\textbf{Proof.}
We give the concentration and approximation steps because the branch properties are conclusions, not assumptions. First,
\[
\mathbb E\widehat G_B(h)=h^{a+1}\int_0^2u^a\mathbf R_1(u)\mathbf R_1(u)^\top g_S(hu)\,du
\succeq c_{S,-}h^{a+1}H_{a,p},
\]
and
\[
\mathbb E\widehat G_I(h,J)=\int_h^1\mathbf b_{h,J}(x)\mathbf b_{h,J}(x)^\top g_S(x)\,dx
\succeq c_{S,-}L_{h,J}.
\]
The boundary summands have operator envelope \(C\), before the common scale \(h^{a+1}\), and effective occupancy \(nh^{a+1}\). For the interior matrix, bounded overlap of the normalized B-splines gives
\[
\left\|\mathbf 1\{X_i\ge h\}X_i^{-a}\mathbf b_{h,J}(X_i)\mathbf b_{h,J}(X_i)^\top\right\|_{\mathrm{op}}\le Ch^{-a},
\]
while the variance proxy is at most \(Ch^{-a}J^{-1}\). Applying the bounded-summand matrix Bernstein lemma at deviation \(cJ^{-1}\), and using the B-spline Gram eigenvalue bound, gives
\[
P\{(\mathcal E_n^G)^c\}\le C J\exp\{-cnh^a/J\}+C\exp\{-cnh^{a+1}\}.
\tag{1}
\]
This is the required actual summand envelope; no boundary-count heuristic is substituted for it.

The same Bernstein calculation for the target loading vectors, followed by the two inverse-Gram bounds on \(\mathcal E_n^G\), yields
\[
\sum_i\Gamma_i^2
=\frac1n\widehat l_B^\top\widehat G_B^{-1}\widehat Q_B\widehat G_B^{-1}\widehat l_B
+\frac1n\widehat l_I^\top\widehat G_I^{-1}\widehat Q_I\widehat G_I^{-1}\widehat l_I
+2\mathcal C_{BI}
\le \frac{CV_a(h)}n,
\tag{2}
\]
where \(\widehat Q_B=P_n[\mathbf1\{X<2h\}\mathbf R_h\mathbf R_h^\top]\), \(\widehat Q_I=P_n[\mathbf1\{X\ge h\}X^{-2a}\mathbf b\mathbf b^\top]\), and \(\mathcal C_{BI}\) is the overlap cross-term on \([h,2h]\). Cauchy--Schwarz absorbs \(\mathcal C_{BI}\). Integrating the envelope in the interior gives
\[
\int_h^1x^{-a}\,dx\asymp V_a(h),
\]
and the boundary quadratic form is \(O\{h^{1-a}/n\}\), which is never larger than \(CV_a(h)/n\). The corresponding first-moment quadratic forms give \(\sum_i|\Gamma_i|\le C\). At a single observation the same inverse bounds give
\[
\frac{\max_i\Gamma_i^2}{\sum_i\Gamma_i^2}
\le C\left\{(nh^{a+1})^{-1}+Jh^{-a}/n\right\}.
\tag{3}
\]
The product-smoothness inequality is equivalent to \(st>1/4\). Substitution of \(h_n\) and \(J_n=\lceil\rho_n^{-1/(s+t)}\rceil\) into (1)--(3), together with \(m/n\to\kappa\), shows that \(nh_n^{a+1}\) and \(nh_n^a/J_n\) dominate \(\log^2(en)\), and the right side of (3) is \(o\{1/\log(en)\}\). Enlarging the fixed certificate constant if necessary therefore proves the asserted failure probability and all coefficient conditions.

It remains to bound the noise-free design discrepancy. On \([0,h]\), Taylor's formula for the order-\(p\) boundary polynomial gives a target integral remainder \(Ch^{s+1}\); multiplying the source remainder \(Ch^s\) by the boundary loading first moment \(O(h)\) gives the same order. On \([h,1]\), insert the population \(L^2(g_S)\) projection. The deterministic projection bias is \(CJ^{-(s+t)}\) by the projection-product lemma. The target empirical error has root mean square at most \(Cm^{-1/2}\). The centered source design term has conditional quadratic scale bounded by (2), and the bounded projection remainder only reduces that bound, giving \(Cn^{-1/2}V_a(h)^{1/2}\). Empirical-versus-population Gram perturbations are controlled by (1) and the resolvent identity
\[
\widehat G^{-1}-G^{-1}=\widehat G^{-1}(G-\widehat G)G^{-1};
\]
their root mean square is bounded by the same two stochastic terms. Adding the boundary and interior pieces proves the last display. \(\square\)